% Generated draft: requires mathematical and source review.
\documentclass[11pt]{article}
\usepackage[margin=25mm]{geometry}
\usepackage{fontspec}
\setmainfont{Latin Modern Roman}
\usepackage{amsmath,amssymb,mathtools,unicode-math}
\setmathfont{Latin Modern Math}
\usepackage{enumitem,needspace,xurl,hyperref}
\setlength{\parindent}{0pt}
\setlength{\parskip}{5pt}
\setlength{\emergencystretch}{3em}
\newcommand{\N}{\mathbb{N}}
\newcommand{\Z}{\mathbb{Z}}
\newcommand{\Q}{\mathbb{Q}}
\newcommand{\R}{\mathbb{R}}
\newcommand{\C}{\mathbb{C}}
\newcommand{\F}{\mathbb{F}}
\newcommand{\A}{\mathbb{A}}
\newcommand{\PP}{\mathbb{P}}
\newcommand{\E}{\mathbb{E}}
\newcommand{\eps}{\varepsilon}
\newcommand{\dd}{\,\mathrm{d}}
\begin{document}
\section*{UnsolvedMath ID: 30002593}
\textbf{Source problem code:} catalogue entry
\subsection{Definitions and mathematical statement}
For a positive integer \(n\), the Artin braid group \(B_n\) is the group of braids on \(n\) strands, with multiplication given by concatenation. For \(n\geq 2\), its standard presentation is
\[
B_n=\left\langle \sigma_1,\ldots,\sigma_{n-1}\ \middle|\ \begin{aligned}
\sigma_i\sigma_j&=\sigma_j\sigma_i &&(|i-j|\geq 2),\\
\sigma_i\sigma_{i+1}\sigma_i&=\sigma_{i+1}\sigma_i\sigma_{i+1} &&(1\leq i\leq n-2)
\end{aligned}\right\rangle.
\]
Under the usual convention, \(B_1\) is trivial. The supplied record does not specify the intended lower bound on the braid index.

The source denotes deformation \(K\)-theory by \(K_{\mathrm{def}}\), a space of irreducible representations by \(\operatorname{Irr}\), and homology by \(H_*\). It does not specify enough conventions to define these expressions uniquely in this question: the representation category, topology on the irreducible-representation space, homology coefficients, and spectrum or graded-group interpretation of deformation \(K\)-theory are unresolved.

\emph{Questions as supplied.} For Artin braid groups \(B_n\), do the following identifications hold generally in the braid index?
\begin{enumerate}
\item Does one have
\[
K_{\mathrm{def}}(B_n)\cong H_*(\operatorname{Irr}(B_n))?
\]
\item Is the deformation \(K\)-theory of \(B_n\) isomorphic to its complex topological \(K\)-theory?
\end{enumerate}
The first display preserves the supplied formula; it does not assert that the two sides have already been specified as objects of the same category. The coefficient system, grading, and meaning of the isomorphism must be supplied before it becomes a precise comparison. For the second question, the underlying space or algebra whose complex topological \(K\)-theory is intended, the variance and grading, and the proposed comparison are likewise unspecified. In particular, interpreting this phrase as the complex \(K\)-theory of the classifying space of \(B_n\) would require an additional source clarification.

The original wording asks whether an unspecified deformation \(K\)-theory theorem extends to Artin braid groups. The record does not state that theorem or its hypotheses, so its relationship to the first displayed formula cannot be fixed from the supplied evidence. These unresolved qualifications prevent a unique fully specified formulation.
\subsection{Short English statement}
The supplied questions ask whether deformation \(K\)-theory of Artin braid groups is described by homology of irreducible representations and whether it agrees with complex topological \(K\)-theory. The relevant representation space, homology coefficients, gradings, comparison target, and referenced extension theorem need clarification. This is a statement-only draft requiring human mathematical review.
\subsection{Sources}
\begin{itemize}
\item Supplied local UnsolvedMath registry record, permanent ID \(30002593\), source problem code \texttt{catalogue entry}, title: \emph{Deformation K-Theory of Artin Braid Groups}; collection: \emph{Oberwolfach Reports Open Problems}. No individual proposer is supplied; the proposed year is recorded as 2014.
\item Supplied report citation: \emph{Algebraic Structures in Low-Dimensional Topology} (2015), Oberwolfach Reports, pp.~1403--1458, DOI \texttt{10.4171/owr/2014/26}. Supplied source URL and literature-source link: \url{https://doi.org/10.4171/owr/2014/26}. The literature-source label attributes this reference to Bryden. The record describes Oberwolfach Reports as documenting mathematical workshops and their open-problem discussions. The citation's 2015 label and its 2014 DOI component are retained as supplied.
\item Supplied literature reference: Azzali et al., \emph{K-theory of pure braid groups}, \url{https://arxiv.org/abs/2105.07848}.
\item Literature assessment dated exactly \texttt{2026-08-21}: the supplied curation labels the two Artin-braid questions open and states that later work computes selected ordinary \(K\)-theory groups but does not establish the requested general representation-ring construction or the deformation-\(K\)-theory isomorphism. This is supplied dated curation, described as literature triage derived from the OWR curation; its status and cited claims have not been independently verified here.
\item The registry assigns the statement status \texttt{corrected\_verified} and reports that it was checked against the official Oberwolfach report, with the clean statement restoring the source's hypotheses and notation. This is a supplied verification claim, not an independent verification in this draft. The full report passage and the referenced deformation \(K\)-theory theorem are not included in the supplied record.
\end{itemize}
\end{document}
